\section{Introducción}
\label{sec:introduccion}

La clasificación manual de piezas pequeñas (brocas, escariadores,
tornillos) en procesos de ferretería o manufactura es una tarea
repetitiva propensa a error humano. PARTS\_SORTING\_PLATFORM busca
automatizar la identificación de estas piezas a partir de una
fotografía, dejando la decisión final y el registro en manos de un
operario, sin incorporar actuadores de separación física de piezas.

El proyecto contempla dos escenarios de uso con requisitos distintos:

\begin{itemize}
  \item Un escenario \textbf{en la nube}, donde varios operarios desde
  distintos dispositivos móviles suben fotografías autenticadas a una
  infraestructura administrada (Railway), reciben el resultado en la
  aplicación y, opcionalmente, un dispositivo ESP32 emparejado muestra
  el resultado en tiempo real.

  \item Un escenario \textbf{local}, pensado para un banco de pruebas o
  una banda transportadora física sin acceso a internet: un teléfono
  Android crea una red WiFi propia (hotspot) a la que se conectan tanto
  el microcontrolador ESP32 como la computadora que ejecuta el modelo
  de clasificación.
\end{itemize}

Ambos modos comparten el mismo propósito funcional (identificar la
pieza y, en el modo local, además estimar su diámetro) pero no
comparten infraestructura ni protocolo de transporte entre sí.
